\documentclass[11pt,a4paper]{article}
\usepackage[T1]{fontenc}
\usepackage[utf8]{inputenc}
\usepackage{amsmath,amssymb,amsthm}
\usepackage[margin=1in]{geometry}
\usepackage[colorlinks=true,linkcolor=blue,urlcolor=blue]{hyperref}
\theoremstyle{plain}
\newtheorem{theorem}{Theorem}
\newtheorem{lemma}[theorem]{Lemma}
\newtheorem{proposition}[theorem]{Proposition}
\newtheorem{corollary}[theorem]{Corollary}
\theoremstyle{definition}
\newtheorem{remark}[theorem]{Remark}

\title{The empirical column recurrences for the table OEIS A263964}
\author{Adrian Perez Fontelles\\ \small Independent researcher}
\date{1 September 2026}
\begin{document}
\maketitle

\begin{abstract}
OEIS A263964 is a two-dimensional table $T(n,k)$, and it carries a block of empirical
recurrences contributed by R.~H.~Hardin --- one for each of its first few columns. They are
true. Column $k$ of the table is an ordinary fixed-width array count, so its rows are the
vertices of a finite digraph and $T(n,k)$ is a number of walks in it; being a walk count the
column is $C$-finite, and whether a given recurrence annihilates it is decided by a finite
exact computation, with the Cayley--Hamilton theorem bounding the work. This note settles the columns $k=1, 2$.
\end{abstract}

\noindent\small 2020 Mathematics Subject Classification. 05A15, 05B45, 15B36.\normalsize

\section{The table and the conjectures}

OEIS A263964 is ``T(n,k)=Number of (n+1)X(k+1) arrays of permutations of 0..(n+1)*(k+1)-1 with each element having index change (+-,+-) 0,0 1,1 or 1,2.''. It has offset $1$, is
read by antidiagonals, and begins
\[
4, 20, 9, 108, 141, 25, 465, 2089,\ \dots
\]
The entry carries the following block:
\begin{quote}\small
Empirical for column k:\\
k=1: a(n) = 2*a(n-1) +2*a(n-2) -a(n-3)\\
k=2: a(n) = 8*a(n-1) +40*a(n-2) -119*a(n-3) -124*a(n-4) +240*a(n-5)\\
k=3: [order 15]\\
k=4: [order 82]\\
\end{quote}
As of the ``Last modified'' line on the live entry (Nov 10 2015 10:24:48, revision 6) these are still
recorded as empirical, and nothing on the entry records any of them as settled.

\section{A column of the table is a walk count}

Fix $k$. The entry defines $T(n,k)$ as the number of arrays of a shape in which one dimension
is $n$ and the other is determined by $k$, subject to a condition imposed on every
$2\times2$ subblock --- a condition local to two consecutive rows and two consecutive columns.
Substituting the fixed value of $k$ into the entry's own wording turns the definition into an
ordinary one-parameter array count, and for such a count the rows are the vertices of a finite
digraph $G_k$: $r\to s$ is an edge exactly when placing $s$ under $r$ satisfies the condition
in every column pair. An admissible array is then a walk, so with $M_k$ the adjacency matrix
of $G_k$,
\[
T(n,k)\;=\;c_k\,\mathbf{1}^{\!\top}M_k^{\,n}\mathbf{1}
\]
for the constant $c_k$ that the entry's name supplies (a factor such as $\tfrac12$ or
$\tfrac14$ where the name says so, and $1$ otherwise).

\begin{lemma}\label{lem:crit}
Let $q(t)=t^{r}-\sum_i c_it^{r-i}$ be the characteristic polynomial of a conjectured
recurrence of order $r$ for column $k$. Then for every $n\ge r$,
\[
T(n,k)-\sum_i c_i\,T(n-i,k)\;=\;c_k\,\mathbf{1}^{\!\top}M_k^{\,n-r}q(M_k)\mathbf{1},
\]
so the recurrence holds from some point on if and only if
$\mathbf{1}^{\!\top}M_k^{j}q(M_k)\mathbf{1}=0$ for every $j\ge0$; and by the
Cayley--Hamilton theorem it is enough to check $j<S_k$, where $S_k$ is the number of vertices
of $G_k$.
\end{lemma}

\begin{proof}
Substitute the walk expression and factor $M_k^{n-r}$ out of $M_k^{n}-\sum_ic_iM_k^{n-i}$;
the scalar $c_k$ is nonzero. For the bound, $M_k^{S_k}$ is an integer combination of
$I,M_k,\dots,M_k^{S_k-1}$, so each $\mathbf{1}^{\!\top}M_k^{j}q(M_k)\mathbf{1}$ with
$j\ge S_k$ repeats that combination of earlier values.
\end{proof}

\section{The computation, column by column}

\begin{center}
\begin{tabular}{ccccc}
line & order & vertices $S$ & proved for & terms checked \\ \hline
column $k$$=1$ & $3$ & $16$ & $>1$ & $7$ \\
column $k$$=2$ & $5$ & $64$ & $>3$ & $8$ \\
\end{tabular}
\end{center}

In each case the vector $w=q(M_k)\mathbf{1}$ was formed by $r$ matrix--vector products in
exact integer arithmetic and $\mathbf{1}^{\!\top}M_k^{j}w$ evaluated for $j=0,1,\dots$,
again exactly; every one of those integers vanishes from the stated point onwards and the run
continued until the working vector was identically zero, which settles all larger $j$ at once.

\subsection*{Column $k$$=1$}
The sequence is ``Number of (n+1)X(2) arrays of permutations of 0..(n+1)*(2)-1 with each element having index change (+-,+-) 0,0 1,1 or 1,2.''. Its rows are the vertices of a digraph on
$S=16$ vertices, edges being the admissible consecutive pairs, and the count is
$\mathbf{1}^{\!\top}M^{n}\mathbf{1}$ up to the constant factor in the entry's name.
The conjectured recurrence
\[
a(n)\;=\;2\,a(n-1) + 2\,a(n-2) - a(n-3)
\]
holds from index $1+1$ on.

\subsection*{Column $k$$=2$}
The sequence is ``Number of (n+1)X(3) arrays of permutations of 0..(n+1)*(3)-1 with each element having index change (+-,+-) 0,0 1,1 or 1,2.''. Its rows are the vertices of a digraph on
$S=64$ vertices, edges being the admissible consecutive pairs, and the count is
$\mathbf{1}^{\!\top}M^{n}\mathbf{1}$ up to the constant factor in the entry's name.
The conjectured recurrence
\[
a(n)\;=\;8\,a(n-1) + 40\,a(n-2) - 119\,a(n-3) - 124\,a(n-4) + 240\,a(n-5)
\]
holds from index $3+1$ on.


\section{Lines whose recurrence the entry does not print}

For the lines below the entry records only an order, the coefficients being in a linked file.
They can still be settled, and the recurrences recovered. A line is a walk count on $S$
vertices, so it satisfies some monic recurrence of order at most $S$ and Berlekamp--Massey on
$2S$ exact terms returns the MINIMAL one. Where that order equals the stated order --- and
where the same holds after the first few terms are dropped, since the entry's recurrence may
begin only partway along --- any recurrence of that order the line satisfies has a
characteristic polynomial that is a multiple of the minimal one and of equal degree, hence is
that polynomial. So the entry's recurrence is the one displayed, whatever its file contains.

\subsection*{Column $k=3$, recovered}
The entry gives only the order for this line: ``k=3: [order 15]''. Its model is a walk on
$S=256$ vertices, so it satisfies a monic recurrence of order at most $S$;
Berlekamp--Massey on $2S=512$ exact terms returns the minimal one, of order
$15$ --- the stated order --- with integer coefficients
\begin{center}\small\begin{tabular}{cccccccc}
$c_{1}$ & $59$ & $c_{5}$ & $-1193960$ & $c_{9}$ & $-3955605504$ & $c_{13}$ & $461708984320$ \\
$c_{2}$ & $-800$ & $c_{6}$ & $-8299248$ & $c_{10}$ & $24437587968$ & $c_{14}$ & $-618475290624$ \\
$c_{3}$ & $-5123$ & $c_{7}$ & $121460288$ & $c_{11}$ & $-38600179712$ & $c_{15}$ & $274877906944$ \\
$c_{4}$ & $202521$ & $c_{8}$ & $-182636544$ & $c_{12}$ & $-99925098496$ &  &  \\
\end{tabular}\end{center}
Repeating with the first $15+4$ terms removed gives order $15$ again.

\subsection*{Column $k=4$, recovered}
The entry gives only the order for this line: ``k=4: [order 82]''. Its model is a walk on
$S=1024$ vertices, so it satisfies a monic recurrence of order at most $S$;
Berlekamp--Massey on $2S=2048$ exact terms returns the minimal one, of order
$82$ --- the stated order --- with integer coefficients
\[
(c_1,\dots,c_{82})=(118,\ 2127,\ -251434,\ 2418222,\ 97287877,\ -1281284386,\ -15907119343,\ 321677121988,\ -160843061642,\ -30899161556508,\ -37928125832132,\ 3872778546393564,\ 11376859061200478,\ -319358239858596494,\ -894710161778357132,\ 15987539408815649436,\ -15987781479726888708,\ -449462363781438481474,\ 6055284735394485056038,\ -9755626948388894497931,\ -233997516144576625400834,\ 969147169186462184804253,\ 970231428358621403848596,\ -21853087189323512847514088,\ 96780840952925083073577519,\ 121999674230571032464701622,\ -1802772104960969825859714805,\ 2743513216505937329805713142,\ 9499085879968597715184654668,\ -36106031314992565338844415104,\ 47881300498791675701796554528,\ 35863863951559978602886503808,\ -602485831863021293257504283136,\ 944990456239565043715611830144,\ 559873354531171697749646659072,\ -2042846405002144061780446031872,\ 5639710350355334732794541880320,\ -10919566159120997303140701454336,\ 18147323560883260096416644837376,\ 39936529667185911745985162797056,\ -114956946568175207078313684893696,\ -182219799022153744081064713322496,\ -155804022154200802577955105996800,\ 1882548460953297450194617397805056,\ -285850169828023917305160597766144,\ -9458814330023695915493376085458944,\ 10611463280095735779833770841997312,\ 26858845523733132505156327751811072,\ -45351653659540425493462454462578688,\ -51721893208246634620152119065313280,\ 128153962852738923260061818732675072,\ 72309172290880790264088521468805120,\ -287601760268876536722533839324315648,\ -81748882298954611457217204880146432,\ 412991900758674169256073756151906304,\ 94618822951706107415260712822374400,\ -310127540725702211914273444128096256,\ -55673602238528243177149497188810752,\ 99063598679677544637623177565962240,\ -112885051673054706587921553354129408,\ -32118319565038796223590575240118272,\ 219629632788546769870763060210171904,\ 67801165499428846487333764314693632,\ -136017390279150443251682106012598272,\ -48912840378723585749768548625940480,\ 29599986855090693656958900142866432,\ 14058551158827272275321597481451520,\ 1996510268761030655279491855679488,\ -6376335876703166987051390600216576,\ -1543477766131110013355911106527232,\ 2928929832630939960301953751187456,\ -141938317623710974670626411773952,\ -343386812814652966586013047586816,\ 188061853093331689484504428183552,\ 9742318402429894060485134778368,\ -13934488740089651776734818729984,\ 1234529994836267689742088273920,\ 74517394163476620228492787712,\ 29128614276432134478604795904,\ -26207698380515738563867836416,\ 773712524553362671811952640,\ 43521329506126650289422336).
\]
Repeating with the first $82+4$ terms removed gives order $82$ again.

\section{Verification}

Three checks, each independent of the algebra above.

First, each model was compared against the entry's own published values for its line, taken
both from the antidiagonal DATA read in either orientation and from the ``Table starts''
block, and agrees at every available term. This is what ties a model to the table:
substituting the fixed index into the entry's wording is a reading of English, and a wrong
reading gives a different digraph and different counts. Across the lines settled here the
model reproduces $15$ published values, $60$ decimal digits in all, every one of
them exactly; a misreading of the entry would have to reproduce all of them by accident.

Second, each conjectured recurrence was evaluated on those published values in exact
integer arithmetic, with no matrices involved, and holds wherever the proved range and the
data overlap.

Third, the annihilation test of Lemma~\ref{lem:crit} was carried out in exact integer
arithmetic rather than on a sampled prefix.

\begin{remark}
A row of the table is a column of the transposed table: fixing $n$ and letting $k$ vary gives a
count of arrays of fixed height, which is the same kind of object with the two dimensions
exchanged. Both are settled here by the same argument. Only the lines the entry states
explicitly, and for which enough published values exist to run the first check above, are
claimed.
\end{remark}

\begin{thebibliography}{9}
\bibitem{oeis} The OEIS Foundation, \emph{The On-Line Encyclopedia of Integer
Sequences}, \url{https://oeis.org}, 2026. Sequence A263964.
\bibitem{stanley} R.~P.~Stanley, \emph{Enumerative Combinatorics}, Volume 1, 2nd ed.,
Cambridge University Press, 2011. (Transfer-matrix method, Section 4.7.)
\bibitem{fs} P.~Flajolet and R.~Sedgewick, \emph{Analytic Combinatorics}, Cambridge
University Press, 2009.
\bibitem{horn} R.~A.~Horn and C.~R.~Johnson, \emph{Matrix Analysis}, 2nd ed., Cambridge
University Press, 2013.
\end{thebibliography}

\end{document}
